\section{A stopped-service common scalar for the final rank-two
seven}

\subsection{1. Scope and audit status}

This note gives the independently audited recurrence proof for exactly
the seven pairs selected by \emph{rank\_\texttt{\seqsplit{two\_mixed\_profile\_7.py}}}. It
composes:

\begin{enumerate}
\def\labelenumi{\arabic{enumi}.}
\tightlist
\item
  a localized Perron--Frobenius activation trial at every dormant
  nonservice vertex;
\item
  a full-chain density-dependent service window after activation; and
\item
  one common scalar on the stopped wedges and their pointwise-good
  complement.
\end{enumerate}

The exact seven-pair fingerprint is

\begin{CodeBlock}
93717536ce82eceefe6909c62568afab31e06695dada8b69defb93335d576957
\end{CodeBlock}

The final proof replaces an invalid geometric-tail shortcut for the
total seed count, distinguishes every birth and death count in the
episode, uses random-time-change exponential estimates for the service
clock, and restores all physical reactions through a deterministic
endpoint envelope. Independent replay passed the stopping, endpoint,
ODE, moment, gluing, and arbitrary-orientation obligations. The exact
seven-pair analytic and recurrence flags are true; global T3-2 remains
false.

\subsection{2. Exact geometry and the common
scalar}

Every pair has lower linkage

\[
                         0\rightleftarrows C             \tag{2.1}
\]

and a strongly connected top linkage of molecularity two and
stoichiometric rank two. The exact service-zero geometry is

\[
\begin{array}{c|c|c}
T&\text{sources on }C=0&\text{dormant nonservice vertices}\\ \hline
\{2A,2C,AB,AC,BC\}&\{2A,AB\}&B\\
\{2A,2C,AB,BC\}&\{2A,AB\}&B\\
\{2A,AB,AC,BC\}&\{2A,AB\}&B\\
\{2B,2C,AB,AC\}&\{2B,AB\}&A\\
\{2B,2C,AB,AC,BC\}&\{2B,AB\}&A\\
\{2B,AB,AC,BC\}&\{2B,AB\}&A\\
\{2C,AB,AC,BC\}&\{AB\}&A,B.
\end{array}                                             \tag{2.2}
\]

There are eight dormant rows. Every pair contains \(AB\), has no
two-active failed descriptor, and has one all-active failed descriptor
with workload \((1,1,1)\).

Put

\[
\begin{split}
 n&=A+B+C,\qquad h=1+n,\\
 F(x)&=K_F+\sum_{i=A,B,C}\log(x_i!),\qquad G=1+F,\\
 V(x)&=G(x)^4+\lambda h(x)^6,\qquad \lambda>0,
\end{split}                                             \tag{2.3}
\]

where \(K_F\) makes \(G\ge1\). This is the single proper scalar used at
every gluing endpoint.

\subsection{3. Pointwise drift away from dormant
wedges}

Write \(\alpha>0\) and \(\delta>0\) for the rates of \(0\to C\) and
\(C\to0\). The top linkage preserves \(n\), and hence

\[
\begin{split}
 {\cal L}h^6
 &=\alpha\{(h+1)^6-h^6\}
   +\delta C\{(h-1)^6-h^6\}\\
 &\le K_0h^5-c_0Ch^5.
\end{split}                                             \tag{3.1}
\]

At a dormant vertex \(d\), relabel \(d\) as \(X\), the service species
\(C\) as \(A\), and the remaining species as \(B\). Section 4 constructs
\(R_d=v_AA+v_BB\), with positive coefficients. Choose a common
sufficiently small \(\eta\), after the fixed pair, orientation, and
rates are given, and define

\[
 {\cal W}_d=\{R_d\le\eta n\},\qquad
 {\cal W}=\bigcup_{d\in D(\mathrm{pair})}{\cal W}_d.    \tag{3.2}
\]

The one or two wedges are disjoint above a finite population sublevel.
Every divergent sequence outside \({\cal W}\) has a subsequence on which

\[
                         {\cal L}V\longrightarrow-\infty. \tag{3.3}
\]

The exact-tier cases are as follows.

\subsubsection{3.1 All-active failure}

Here \(A,B,C=\Theta(N)\), \(h=\Theta(N)\), and the top source rates are
\(\Theta(N^2)\). The exact discrete fourth-power estimate from the
audited homogeneous rank-two scalar theorem gives

\[
                    ({\cal L}G^4)^+
                    \le K N^5(\log N)^3.                \tag{3.4}
\]

The service term in (3.1) has negative order \(N^6\), which is a strict
polynomial margin.

\subsubsection{\texorpdfstring{3.2 Passing cones with
\(C\to\infty\)}{3.2 Passing cones with C\textbackslash to\textbackslash infty}}

The powered exact-tier estimate gives

\[
 {\cal L}G^4\le-cG^3A_Ng_N,\qquad g_N\longrightarrow\infty,             \tag{3.5}
\]

after its lower-order binomial terms are absorbed. Here \(A_N\) is the
largest enabled source propensity. At the same time, (3.1) is eventually
negative. Both summands of \(V\) are generator-good.

\subsubsection{\texorpdfstring{3.3 Passing cones with \(C\)
bounded}{3.3 Passing cones with C bounded}}

Now \(h\asymp A+B\). On the service-zero edge, (2.2) gives the coercive
source explicitly. A non-dormant pure endpoint has a \(2A\) or \(2B\)
source of order \(h^2\). If both \(A\) and \(B\) have positive limiting
fractions, the universal \(AB\) source has order \(h^2\). The only
omitted pure endpoints are the wedges in (3.2). Compactness of their
complement therefore gives

\[
                              A_N\ge c h^2.             \tag{3.6}
\]

Since \(G\asymp h\log h\), the negative magnitude in (3.5) is at least

\[
                         ch^5(\log h)^3g_N,             \tag{3.7}
\]

which absorbs the \(O(h^5)\) birth term in (3.1).

There are no two-active failures. Every one-active failure enters one of
the dormant wedges, and an affine-infeasible descriptor cannot occur in
a fixed population class. The bad-sequence contradiction turns (3.3)
into

\[
                 {\cal L}V\le-1                         \tag{3.8}
\]

outside \({\cal W}\) and one finite set.

\subsection{4. Localized PF trials and the exact activation
counts}

Fix a dormant chart. Its normalized top support is one of

\[
\{2A,2B,AB,AX,BX\},\quad
\{2A,2B,AX,BX\},\quad
\{2A,AB,AX,BX\},\quad
\{2B,AB,AX,BX\}.                                      \tag{4.1}
\]

Killing each carrier edge which enters an inactive quadratic gives a
transient two-state generator \(Q\). Let

\[
 q=(-Q)^{-1}{\bf1}>0,\qquad v_i=1-\varepsilon q_i>0.   \tag{4.2}
\]

For sufficiently small \(\varepsilon\), the rate-weighted reward at both
carrier sources is positive. Thus, exactly for falling-factorial
propensities,

\[
 R=v_AA+v_BB,\qquad
 {\cal L}_TR\ge cXR-KR^2.                              \tag{4.3}
\]

The physical birth increases \(R\); the death contributes at worst
\(-KR\). After reducing \(\eta\), for \(0<R\le\eta n\) and all large
\(n\),

\[
 {\cal L}R\ge c_1nR,\qquad
 \Gamma R:=\sum_r\lambda_r(\Delta_rR)^2\le C_1nR.      \tag{4.4}
\]

All jumps of \(R\) are bounded. For a sufficiently small \(\theta>0\),

\[
 {\cal L}e^{-\theta R}
 \le-c_2\theta nR e^{-\theta R}.                       \tag{4.5}
\]

Starting with one physical \(0\to C\) seed and stopping at \(R=0\) or
\(R\ge\eta n\), optional stopping gives a success probability

\[
 p_d\ge1-e^{-\theta v_A}+o(1)>0.                       \tag{4.6}
\]

This probability statement alone does not imply a geometric tail for the
total number of births. The required repair is a localized trial
estimate. Choose \(r_0\) larger than twice the maximum \(R\)-jump. On
\(r_0\le R\le\eta n\), Taylor's formula yields

\[
                         {\cal L}\log R\ge c_3n.        \tag{4.7}
\]

Below \(r_0\), only finitely many inactive phases occur and every
enabled carrier race has order \(n\). The killed-carrier cut excludes a
closed nonabsorbing finite phase. Dynkin applied to (4.7), followed by
the strong Markov property over dyadic bands and the finite lower phase,
gives constants \(a,b,C>0\) such that the fast part \(\rho_n\) of a
trial obeys

\[
 \mathbb P\!\left\{\rho_n>u\,{C\log n\over n}\right\}
 \le ae^{-bu},\qquad u\ge1,                            \tag{4.8}
\]

until activation, return to a dormant vertex, or a population-shell
exit. This is a physical-time estimate, not a fast-jump count.

Let \(J\) be the number of extra \(0\to C\) births in the fast part.
Since \(\rho_n\) depends on the birth clock, \(J\) is not asserted to be
Poisson conditional on \(\rho_n\). Instead, the stopped counting-process
exponential martingale gives

\[
 \mathbb E\exp\{sJ-\alpha(e^s-1)\rho_n\}\le1.           \tag{4.9}
\]

Apply Cauchy--Schwarz to the martingale with parameter \(2s\). Equations
(4.8)--(4.9) then give a uniform exponential moment for \(J\), and that
moment tends to one as \(n\to\infty\). At a dormant vertex, one
additional seed is charged before the next fast trial. The number of
localized trials is dominated by a geometric variable with success
parameter

\[
 p_*=\min_{d\in D(\mathrm{pair})}p_d>0,                \tag{4.10}
\]

where the minimum is taken only over the one or two dormant charts of
the fixed pair after its rates are fixed. A completed macroepisode may
be re-marked at the other dormant vertex; the next application uses the
same finite minimum. No constant is claimed uniformly over rate vectors.

Let \(N_0\) be the population at the beginning of the episode, \(K\) the
total number of activation-stage births, and \(D_{\rm pre}\) the total
number of activation-stage deaths. The compound-geometric calculation
from (4.9)--(4.10), including an arbitrary initial point in a wedge,
gives some \(s_0>0\) with

\[
 \sup_{N_0\ge N_*}\mathbb E e^{s_0K}<\infty,\qquad
 \sup_{N_0\ge N_*}\mathbb E\tau_{\rm act}^{\,m}<\infty
 \quad(m\ge1).                                         \tag{4.11}
\]

Localize the episode at \(n\le N_0/2\) or \(n\ge2N_0\). A downward exit
is favorable. An upward exit requires at least \(N_0\) births, so (4.11)
implies

\[
               \mathbb P\{\text{upward exit}\}\le Ce^{-cN_0}.           \tag{4.12}
\]

A downward exit requires order-\(N_0\) service deaths during a
compound-geometric sum of fast trials. The localized trial-time estimate
(4.8), the death counting-process exponential martingale, and a split at
\(\varepsilon N_0/\log N_0\) trials give the bound. Above the split, the
geometric trial tail is \(e^{-cN_0/\log N_0}\). Below it, the summed
compensator is at most \(C\varepsilon N_0\) outside an exponentially
small trial-duration event; choose \(\varepsilon\) first and apply the
counting martingale. Thus, for every \(m\),

\[
 \mathbb P\{\text{activation-stage downward exit}\}
 =o(N_0^{-m}).                                         \tag{4.13}
\]

This tail estimate is used only to remove the branch from the regular
service calculation. Its actual scalar contribution is favorable and is
retained separately in Section 7.

On the regular activation branch, with activation population \(N_a\),
the exact identity is

\[
                         N_a=N_0+K-D_{\rm pre},
 \qquad {N_0\over2}\le N_a\le2N_0.                    \tag{4.14}
\]

No lower clock has been removed in (4.8)--(4.14).

\subsection{5. Infinite top-ODE service
integral}

On the unit simplex, let \(z=(a,b,c)\) solve

\[
 {dz\over d\tau}
 =\sum_{y\to y'\in T}\kappa_{yy'}z^y(y'-y),
 \qquad a+b+c=1.                                      \tag{5.1}
\]

Let \({\cal A}\) be the finite union of normalized activation shells. It
is compact and separated from every dormant vertex.

\begin{quote}
\textbf{Lemma 5.1.} For every \(z(0)\in{\cal A}\), \[
                      \int_0^\infty c(\tau)\,d\tau=\infty.             \tag{5.2}
\]
\end{quote}

If the integral were finite, boundedness of the polynomial vector field
would make \(c\) uniformly continuous, and Barbalat's argument would
give \(c(\tau)\to0\). The omega-limit would be a nonempty compact
invariant subset of \(c=0\).

The directed cut in (2.2) identifies the largest invariant subset of
that face. If \(a,b>0\), every present service-free source is enabled,
and strong connectivity gives a cut edge to a \(C\)-containing node, so
\(\dot c>0\). At the pure \(A\) vertex, an existing \(2A\) source either
creates \(C\) or first creates \(B\), after which the same cut applies.
Thus the vertex remains in the face only when \(2A\) is absent. The
corresponding statement holds at the pure \(B\) vertex. Hence the
largest invariant subset consists exactly of the dormant vertices in
(2.2).

An omega-limit contained in this finite set is one vertex. In its chart,
the fluid form of (4.3) gives

\[
                             \dot R\ge cR-KR^2.          \tag{5.3}
\]

For small positive \(R\), the right side is positive, so an orbit which
has left the vertex cannot approach it. This contradiction proves (5.2).

For fixed \(T\), the map

\[
                         I_T(z)=\int_0^T c_z(\tau)\,d\tau               \tag{5.4}
\]

is continuous. The increasing open sets \(\{z:I_T(z)>M\}\) cover the
compact set \({\cal A}\) as \(T\to\infty\). Therefore, for every
prescribed \(M\), there is one finite \(T=T(M)\) such that

\[
                         \inf_{z\in{\cal A}}I_T(z)>M.   \tag{5.5}
\]

No stationary or fast-mixing assumption is used.

\subsection{6. Lattice-uniform full-chain
service}

On the regular activation branch, run every physical clock for time
\(T/N_a\). For the fluid estimate, localize at
\(N_t\notin[N_a/2,2N_a]\); the endpoint analysis below retains either
localized exit and does not replace the physical chain. Let

\[
 Z_{N_a}(\tau)
 =N_a^{-1}X(\tau_{\rm act}+\tau/N_a),\qquad 0\le\tau\le T.              \tag{6.1}
\]

Top clocks have rates \(N_a^2\kappa z^y+O(N_a)\), so their compensated
density martingales have quadratic variation \(O(N_a^{-1})\). Let
\(B_{\rm win}\) and \(D_{\rm win}\) be the birth and death counts in
this window. For the unlocalized deterministic horizon, the birth count
is Poisson with mean \(\alpha T/N_a\); the stopped count is bounded by
that same underlying clock. The death compensator is

\[
 \Lambda_{\rm win}
 =\delta\int_0^{T/N_a}C_t\,dt
 =\delta\int_0^T (Z_{N_a}(\tau))_C\,d\tau.             \tag{6.2}
\]

Since deaths only reduce population,

\[
 \Lambda_{\rm win}
 \le\delta T+
 {\delta T\over N_a}B_{\rm win}.                       \tag{6.3}
\]

The exponential counting-process martingale gives, after localization,

\[
 \mathbb E\exp\{sD_{\rm win}-(e^s-1)\Lambda_{\rm win}\}\le1.            \tag{6.4}
\]

Apply Cauchy--Schwarz to (6.4) with parameter \(2s\), and then use
(6.3). Together with the Poisson exponential moment of \(B_{\rm win}\),
this gives uniform exponential moments for both window counts. Thus the
death clock is controlled through its random compensator; it is not
replaced by a fixed pathwise Poisson clock.

Random time changes, high-moment Burkholder estimates, and Gronwall now
give, uniformly over every lattice point whose normalized value is
within \(O(N_a^{-1})\) of \({\cal A}\),

\[
 \mathbb E\sup_{0\le\tau\le T}
 \left\lVert Z_{N_a}(\tau)-z(\tau)\right\rVert^m
 \longrightarrow0                                    \tag{6.5}
\]

for each fixed \(m\). The lower counts have \(O(1)\) moments and
therefore move density by \(o(1)\). By the compensator identity and
(5.5),

\[
 \inf_{\text{activation lattice states}}\mathbb E D_{\rm win}
 \ge\delta\inf_{z\in{\cal A}}I_T(z)-o(1).              \tag{6.6}
\]

Let

\[
                       \mu_K=\sup_{N_0\ge N_*}\mathbb EK<\infty.
\]

Choose \(M\), hence \(T\), so that \(\delta M>\mu_K+4\). Since
\(\mathbb EB_{\rm win}=o(1)\), define

\[
                       \overline\Delta
                       :=K+B_{\rm win}-D_{\rm win}.     \tag{6.7}
\]

Let \({\cal R}_{N_0}\) denote the regular branch on which neither the
activation stage nor the service window exits its displayed population
shell. Upward activation exits have probability \(O(e^{-cN_0})\) by
(4.11)--(4.12), while (4.13) makes the removed activation-down branch
super-polynomially small. Upward service exits require order-\(N_0\)
births and have an exponential bound. A service-window downward exit
requires order-\(N_0\) deaths; (6.3)--(6.4) make both its probability
and its truncated count moments exponentially small. Each downward
branch is also retained as a separate favorable contribution in Section
7; its negative scalar contribution is not used to improve any
conditional count estimate.

On \({\cal R}_{N_0}\), use the event-weighted inequalities

\[
 \mathbb E[K;{\cal R}_{N_0}]\le\mathbb EK\le\mu_K,
 \qquad
 \mathbb E[B_{\rm win};{\cal R}_{N_0}]=o(1).            \tag{6.8}
\]

Uniformity of (6.6), followed by the exponential count bounds for the
removed service-shell event, gives

\[
 \mathbb E[D_{\rm win};{\cal R}_{N_0}]
 \ge\delta M-o(1).                                     \tag{6.9}
\]

No estimate of the form \(\mathbb E[K\mid{\cal R}_{N_0}]\le\mu_K\) is
asserted or needed. Therefore, for all sufficiently large \(N_0\),

\[
 \mathbb E[\overline\Delta;{\cal R}_{N_0}]\le-3,\qquad
 \sup_{N_0}\mathbb E e^{s_1|\overline\Delta|}<\infty    \tag{6.10}
\]

for some \(s_1>0\). The second, unconditional, statement uses (4.11) and
(6.3)--(6.4). The exact physical endpoint identity, with no count
discarded or counted twice, is

\[
\begin{split}
 N_a&=N_0+K-D_{\rm pre},\\
 N_{\rm end}
 &=N_a+B_{\rm win}-D_{\rm win}\\
 &=N_0+K-D_{\rm pre}+B_{\rm win}-D_{\rm win}
 =N_0+\overline\Delta-D_{\rm pre}.
\end{split}                                             \tag{6.11}
\]

In particular,

\[
                         N_{\rm end}\le N_0+\overline\Delta.            \tag{6.12}
\]

For every \(q>8\), (4.11), (6.10), and the shell localization give

\[
\begin{split}
 \sup_{N_0}\mathbb E\tau^q&<\infty,\\
 \mathbb E(1+N_{\rm end})^q&\le C_q(1+N_0)^q,\\
 \sup_{N_0}\mathbb E\{(N_{\rm end}-N_0)^+\}^q&<\infty.
\end{split}                                             \tag{6.13}
\]

These are the endpoint and duration uniform-integrability estimates used
in the scalar calculation.

\subsection{7. Deterministic endpoint envelope and stopped scalar
drift}

Define the monotone deterministic envelope

\[
 {\cal P}(m)
 =\{1+K_F+\log(m!)\}^4+\lambda(1+m)^6,\qquad m\ge0.    \tag{7.1}
\]

For every state \(y\) of total population \(m\),

\[
                              V(y)\le{\cal P}(m).       \tag{7.2}
\]

Also, the exact multinomial inequality gives

\[
 0\le\log(N_0!)-\sum_i\log(x_i!)
 =\log{N_0!\over A!B!C!}\le N_0\log3.                 \tag{7.3}
\]

Consequently,

\[
                     {\cal P}(N_0)-V(x)
                     \le C N_0^4(\log N_0)^3.          \tag{7.4}
\]

On the regular branch, (6.12) and monotonicity of \({\cal P}\) give the
pathwise deterministic inequality

\[
 V(X_\tau)-V(x)
 \le {\cal P}(N_0+\overline\Delta)-{\cal P}(N_0)
    +C N_0^4(\log N_0)^3.                              \tag{7.5}
\]

This does not delete \(D_{\rm pre}\) from the process or compare two
reaction systems. Every pre-service death remains in the exact identity
(6.11); it only improves the total-population envelope (7.5). All top
reactions are included through the factorial maximum in \({\cal P}\).

The exponential moment in (6.10) permits an exact Taylor expansion with
a uniform remainder. The sixth-power part satisfies

\[
 \mathbb E[\{(1+N_0+\overline\Delta)^6-(1+N_0)^6\};
             {\cal R}_{N_0}]
 =6(1+N_0)^5\mathbb E[\overline\Delta;{\cal R}_{N_0}]
   +O(N_0^4)
 \le-cN_0^5.                                          \tag{7.6}
\]

The positive change of the factorial-maximum part is
\(O(N_0^3(\log N_0)^4)\), while (7.4) costs \(O(N_0^4(\log N_0)^3)\).
Both are strictly below (7.6). Hence

\[
 \mathbb E_x[\{V(X_\tau)-V(x)+\tau\};{\cal R}_{N_0}]
 \le-c_1N_0^5                                         \tag{7.7}
\]

for every sufficiently large wedge state. This is an event-weighted
regular contribution, not a conditional expectation.

At an activation-stage downward shell exit, the endpoint population is
at most \(N_0/2+1\). Equations (7.1)--(7.4) give a negative
order-\(N_0^6\) increment. At an upward exit, the first crossing has
population at most \(2N_0+1\), so its positive envelope is \(O(N_0^6)\);
(4.12) makes its expected contribution \(O(N_0^6e^{-cN_0})\).

For a service-window shell exit, first split on \(K\le N_0/4\). Then
\(N_a\le N_0+K\le5N_0/4\), so a downward crossing of \(N_a/2\), combined
with the deterministic envelope, is favorable at order \(N_0^6\), while
an upward crossing of \(2N_a\) requires order-\(N_0\) window births. The
complement \(K>N_0/4\) has exponentially small probability by (4.11),
and its endpoint is charged by the same exponential moment and
deterministic envelope. Equations (6.3)--(6.4) give the corresponding
exponential window tail. Combining all branches yields

\[
                         \mathbb E_x\{V(X_\tau)-V(x)+\tau\}\le-1        \tag{7.8}
\]

outside a finite wedge sublevel.

\subsection{8. Fixed-class gluing and
nonexplosion}

The sole population-increasing clock has constant rate \(\alpha\). Hence
total population is pathwise bounded above by its initial value plus a
rate-\(\alpha\) Poisson process. On every finite population shell, the
top state space and all reaction rates are finite. This proves
nonexplosion and the finite-time moments used in localization.

Use the partition

\[
 E=K\cup(E\setminus{\cal W})\cup({\cal W}\setminus K), \tag{8.1}
\]

where \(K\) contains the finite exceptions in (3.8) and (7.8). Equation
(3.8) is the generator-good hypothesis of the common-entropy
physical-time gluing theorem; (7.8) is its exceptional-episode
hypothesis for the same globally proper \(V\). The stopping rules are
state-based and strong-Markov compatible. The theorem gives finite mean
hitting of \(K\). After restricting to a closed irreducible population
class, the finite mixture of return bounds from \(K\) gives positive
recurrence.

This proves recurrence for the seven pairs, for each fixed strong top
orientation and positive rate vector. The independent audit replayed
Sections 3--8 and found no load-bearing gap.

\subsection{9. Reproduction and independent
audit}

Run

\begin{CodeBlock}
PYTHONPATH=src python3 -B src/rank_two_mixed_profile_7_stopped_service.py
PYTHONPATH=src python3 -B -m unittest \
  tests/test_rank_two_mixed_profile_7_stopped_service.py -v
\end{CodeBlock}

The independent replay verified:

\begin{enumerate}
\def\labelenumi{\arabic{enumi}.}
\tightlist
\item
  localized PF trial survival, duration, and compound-birth exponential
  moments from arbitrary wedge starts;
\item
  exact identities (4.14) and (6.11), including both shell exits;
\item
  the directed-cut/Barbalat proof of infinite integrated service;
\item
  lattice-uniform fluid convergence and random-compensator exponential
  estimates;
\item
  \(q>8\) endpoint and duration uniform integrability;
\item
  the \(C\)-bounded passing estimate; and
\item
  the deterministic envelope and fixed-class gluing.
\end{enumerate}

All seven scoped recurrence flags are promoted. Global T3-2 remains
uncertified because the disjoint hard 333-pair family is still open.
